\section{X-Band Marine Radar Signal Processing}\label{sec:xband}

The effectiveness of radar-guided camera labeling systems depends fundamentally on the quality of radar detections that drive the cross-sensor association process. As discussed in Chapter~\ref{sec:fusion}, fusion architectures ranging from early-stage feature-level integration to late-stage decision fusion all require reliable radar measurements as input. When radar detections contain excessive false alarms or miss true targets due to inadequate signal processing, the resulting camera labels will be corrupted regardless of how sophisticated the fusion algorithm may be. This chapter examines the signal processing techniques that transform raw radar returns into the detection lists used for camera guidance, with particular emphasis on constant false alarm rate (CFAR) detection, sea clutter modeling, and the characteristics of commercial X-band marine radars such as the Simrad Halo~20+.

X-band marine radar systems, operating at frequencies near 9.4~GHz with wavelengths around 3~cm, provide the spatial resolution necessary for detecting small maritime targets at ranges from tens of meters to several nautical miles \cite{Skolnik2008}. The signal processing chain begins with pulse compression to enhance range resolution and signal-to-noise ratio (SNR), proceeds through CFAR detection to maintain acceptable false alarm rates across varying clutter conditions, and concludes with clustering and tracking algorithms that associate detections across scans \cite{Richards2022}. Understanding this processing chain is essential for researchers implementing radar-camera systems, as design choices at each stage directly impact the accuracy and completeness of automated camera labels.

\subsection{Constant False Alarm Rate Detection}

CFAR detection algorithms adaptively adjust detection thresholds based on local estimates of noise or clutter power, maintaining a constant probability of false alarm across spatially varying background conditions \cite{Finn1968}. The fundamental CFAR threshold equation takes the form:
\begin{equation}
T = \alpha \cdot \hat{Z}
\end{equation}
where $T$ is the detection threshold, $\hat{Z}$ is an estimate of the local noise or clutter power, and $\alpha$ is a scaling factor determined by the desired false alarm rate and the assumed clutter distribution. A range-azimuth cell is declared to contain a target if its magnitude exceeds this adaptive threshold.

The cell-averaging CFAR (CA-CFAR) detector, introduced by Finn and Johnson, estimates background power by averaging samples from leading and lagging reference windows surrounding the cell under test (CUT) \cite{Finn1968}. Guard cells immediately adjacent to the CUT are excluded to prevent target energy from contaminating the noise estimate. CA-CFAR achieves optimal performance in homogeneous clutter environments where all reference cells contain independent and identically distributed clutter samples \cite{Gandhi1988}. The CA-CFAR threshold is computed as:
\begin{equation}
T_{\text{CA}} = \alpha_{\text{CA}} \cdot \frac{1}{N} \sum_{i=1}^{N} Z_i
\end{equation}
where $N$ is the total number of reference cells and $Z_i$ represents the magnitude or power of the $i$-th reference cell. The scaling factor $\alpha_{\text{CA}}$ depends on $N$ and the desired probability of false alarm $P_{\text{fa}}$. For Rayleigh-distributed clutter, $\alpha_{\text{CA}} = N \left( P_{\text{fa}}^{-1/N} - 1 \right)$ \cite{Richards2022}.

When multiple targets appear in close proximity, interfering targets may enter the reference windows and inflate the noise estimate, leading to detection losses. The ordered-statistic CFAR (OS-CFAR) addresses this multi-target scenario by sorting the reference cells in ascending order and selecting the $k$-th smallest value as the noise estimate rather than the mean \cite{Rohling1983}. By choosing $k$ appropriately (typically $k \approx 3N/4$ for $N$ reference cells), OS-CFAR ensures that interfering targets are excluded from the threshold calculation. The OS-CFAR threshold is:
\begin{equation}
T_{\text{OS}} = \alpha_{\text{OS}} \cdot Z_{(k)}
\end{equation}
where $Z_{(k)}$ denotes the $k$-th order statistic of the reference window. OS-CFAR sacrifices approximately 1--2~dB of detection sensitivity compared to CA-CFAR in homogeneous environments but provides superior robustness when interfering targets are present \cite{Gandhi1988}.

Clutter edges, such as the boundary between open water and land, create severe performance degradation for both CA-CFAR and OS-CFAR. When the reference window straddles a clutter transition, the noise estimate becomes biased, resulting in either elevated false alarms (if low-clutter cells dominate the estimate in a high-clutter region) or missed detections (if high-clutter cells dominate in a low-clutter region). The smallest-of CFAR (SO-CFAR) combats clutter edges by computing separate CA-CFAR thresholds from the leading and trailing reference windows and selecting the minimum \cite{Trunk1978}:
\begin{equation}
T_{\text{SO}} = \min \left( \alpha_{\text{L}} \sum_{i \in L} Z_i, \, \alpha_{\text{T}} \sum_{i \in T} Z_i \right)
\end{equation}
where $L$ and $T$ denote the leading and trailing reference cells. SO-CFAR prevents threshold inflation in the presence of clutter edges but suffers increased false alarms in homogeneous regions because the minimum statistic is biased downward. The complementary greatest-of CFAR (GO-CFAR) takes the maximum of the two estimates to reduce false alarms at the cost of increased detection losses at clutter edges \cite{Hansen1972}.

The automatic censored cell-averaging CFAR (ACCA-CFAR) attempts to combine the benefits of CA-CFAR and OS-CFAR by iteratively removing reference cells that exceed a censoring threshold, then averaging the remaining cells \cite{Farina1987}. The algorithm begins with all reference cells, computes a preliminary mean, identifies cells exceeding a multiple of that mean, censors those cells, and recomputes the mean from the censored set. This process continues until convergence or a maximum iteration count. ACCA-CFAR performs well in both homogeneous and multi-target scenarios but requires more computation than simpler CFAR variants.

Recent advances in deep learning have introduced neural network-based CFAR detectors that learn adaptive thresholding strategies directly from labeled radar data. Diskin and colleagues proposed CFARnet, a convolutional neural network trained to predict detection thresholds for each range-azimuth cell given a local neighborhood of radar returns \cite{Diskin2024}. CFARnet demonstrated improved detection performance on maritime targets compared to CA-CFAR and OS-CFAR, particularly in heterogeneous clutter conditions. However, the method requires extensive labeled training data and lacks the analytical performance guarantees of classical CFAR algorithms.

Chen et al. developed Radar-PPInet, a deep network that processes plan-position indicator (PPI) radar images to jointly perform clutter suppression, target detection, and track extraction \cite{Chen2021}. The network architecture combines convolutional layers for feature extraction with recurrent layers to exploit temporal correlation across consecutive radar scans. Evaluation on real-world maritime radar data showed superior detection rates compared to conventional CFAR processing, but generalization to new environments remains a challenge without fine-tuning.

Zhang and colleagues introduced YOLO-SWFormer, adapting the YOLO object detection framework with a Swin Transformer backbone for maritime radar target detection \cite{Zhang2023b}. The approach treats radar PPI images as natural images and applies bounding box regression to localize targets. YOLO-SWFormer achieved real-time inference speeds suitable for shipboard deployment but required careful data augmentation to handle the unique characteristics of radar imagery, such as rotational symmetry and range-dependent resolution.

Marine-Faster R-CNN, proposed by Chen et al., extends the Faster R-CNN framework to maritime radar by incorporating domain-specific modifications including a novel anchor design tailored to the aspect ratios of ship returns and a multi-scale feature pyramid to handle targets at varying ranges \cite{Chen2022}. The method demonstrated state-of-the-art detection accuracy on benchmark datasets but, like other learning-based approaches, depends critically on the availability of large annotated radar datasets, which are scarce in the maritime domain.

Classical CFAR algorithms remain the dominant approach in operational marine radar systems due to their computational efficiency, analytical tractability, and well-understood performance characteristics. However, deep learning CFAR methods represent a promising research direction, particularly for scenarios involving complex clutter environments where traditional model-based approaches struggle. For radar-guided camera labeling applications, the choice of CFAR algorithm directly affects the quantity and quality of radar detections available for cross-sensor association.

\subsection{Sea Clutter Statistical Models}

Accurate modeling of sea clutter statistics is essential for designing CFAR detectors with predictable false alarm rates and for evaluating detection performance through simulation. Sea clutter arises from radar backscatter by ocean surface waves, with statistical properties that depend on wind speed, wave height, radar resolution, polarization, incidence angle, and frequency \cite{Ward2006}.

The Rayleigh distribution represents the simplest sea clutter model, applicable to low-resolution radars observing homogeneous sea states with many independent scatterers per resolution cell \cite{Skolnik2008}. Under the Rayleigh model, the clutter amplitude $A$ follows the probability density function:
\begin{equation}
p_A(a) = \frac{a}{\sigma^2} \exp \left( -\frac{a^2}{2\sigma^2} \right), \quad a \geq 0
\end{equation}
where $\sigma^2$ is the mean clutter power. The corresponding power distribution is exponential:
\begin{equation}
p_Z(z) = \frac{1}{\mu} \exp \left( -\frac{z}{\mu} \right), \quad z \geq 0
\end{equation}
with mean $\mu = 2\sigma^2$. Rayleigh clutter is fully characterized by its mean power, simplifying CFAR design but failing to capture the heavy-tailed behavior observed in high-resolution radars and moderate-to-high sea states \cite{Ward2013}.

The K-distribution provides a more realistic model for sea clutter in many practical scenarios by representing the radar return as the product of two independent random processes: Gaussian speckle and a slowly varying gamma-distributed texture component \cite{Ward1981}. The K-distribution amplitude PDF is:
\begin{equation}
p_A(a) = \frac{4}{\Gamma(\nu)} \left( \frac{\nu}{2\mu} \right)^{(\nu+1)/2} a^{\nu} K_{\nu-1} \left( a \sqrt{\frac{2\nu}{\mu}} \right), \quad a \geq 0
\end{equation}
where $\mu$ is the mean power, $\nu$ is the shape parameter controlling the tail heaviness, and $K_{\nu-1}$ denotes the modified Bessel function of the second kind with order $\nu - 1$. Small values of $\nu$ (typically $0.1 < \nu < 2$) correspond to spiky, heavy-tailed clutter observed at high resolution and high sea states, while large $\nu$ approaches the Rayleigh limit. The K-distribution has been extensively validated against real sea clutter measurements and is widely used in maritime radar performance prediction \cite{Watts2022}.

The compound Gaussian framework generalizes the K-distribution by allowing arbitrary texture distributions while maintaining the Gaussian speckle component \cite{Gao2020}. This flexibility enables fitting a broader range of observed clutter statistics but complicates CFAR detector design because optimal threshold scaling factors must be derived for each specific texture model. Shui et al. demonstrated that compound Gaussian models with inverse Gaussian texture provide improved fits to X-band sea clutter in certain conditions compared to the K-distribution \cite{Shui2019}.

The Pareto distribution, also known as the compound K-distribution, describes high-resolution X-band sea clutter and offers the advantage of enabling simpler optimal detector structures compared to the K-distribution \cite{Farshchian2010}. The Pareto amplitude PDF is:
\begin{equation}
p_A(a) = \frac{2\alpha \beta^\alpha}{a^{2\alpha+1}}, \quad a \geq \beta
\end{equation}
where $\alpha$ is the shape parameter and $\beta$ is the scale parameter. Pareto-distributed clutter exhibits a power-law tail, capturing the extreme sea spikes that characterize high-resolution radar observations. Leung and colleagues showed that Pareto-based CFAR detectors outperform K-distribution-based detectors for small target detection in X-band maritime radar data \cite{Leung2022}.

The Weibull distribution represents another popular model for sea clutter, offering mathematical tractability and the ability to fit a wide range of empirical clutter statistics through its shape parameter $c$ \cite{Wang2024}. The Weibull amplitude PDF is:
\begin{equation}
p_A(a) = \frac{c}{\lambda} \left( \frac{a}{\lambda} \right)^{c-1} \exp \left( -\left( \frac{a}{\lambda} \right)^c \right), \quad a \geq 0
\end{equation}
where $\lambda$ is the scale parameter. The shape parameter controls tail behavior: $c = 2$ yields the Rayleigh distribution, $c < 2$ produces heavy tails, and $c > 2$ gives light tails. Weibull CFAR detectors have been implemented in operational systems, though the Weibull model generally provides poorer fits to high-resolution sea clutter compared to the K-distribution or Pareto models \cite{Mezache2015}.

Recent research has focused on non-stationary and multi-scale clutter models to capture temporal and spatial variations in sea clutter characteristics. Liu et al. proposed a time-varying K-distribution model with slowly evolving shape and scale parameters estimated via Kalman filtering \cite{Liu2023}. Zhang and colleagues developed a spatially heterogeneous clutter model that partitions the radar field of view into regions with distinct statistical properties, enabling adaptive CFAR processing tailored to local clutter conditions \cite{Zhang2023}.

For radar-guided camera labeling systems, the choice of clutter model affects detection performance and false alarm rates. Systems operating in benign conditions (low sea states, long ranges, coarse resolution) may perform adequately with Rayleigh-based CFAR, while those deployed in challenging environments (high sea states, short ranges, fine resolution) require K-distribution, Pareto, or compound Gaussian models to maintain acceptable false alarm rates. Empirical clutter characterization using data collected from the specific radar and environment of interest remains the most reliable approach for selecting an appropriate model.

\subsection{Pulse Compression}

Pulse compression techniques enable marine radars to achieve fine range resolution while maintaining sufficient energy for long-range detection \cite{Richards2022}. A simple pulsed radar transmits a short-duration, high-power pulse to maximize range resolution, but the achievable pulse duration is limited by peak power constraints and transmitter technology. Pulse compression circumvents this limitation by transmitting a longer, frequency-modulated or phase-coded waveform that is compressed upon reception through matched filtering, achieving the range resolution of a short pulse while delivering the energy of a long pulse.

The most common pulse compression waveform for marine radar is the linear frequency modulated (LFM) chirp, in which the instantaneous frequency increases or decreases linearly over the pulse duration $T_p$ \cite{Cook1960}. The transmitted signal is:
\begin{equation}
s(t) = \text{rect}\left( \frac{t}{T_p} \right) \exp \left( j 2\pi \left( f_c t + \frac{B}{2T_p} t^2 \right) \right)
\end{equation}
where $f_c$ is the carrier frequency, $B$ is the bandwidth, and $\text{rect}(t)$ is the rectangular window function. Upon reception, the echo is passed through a matched filter with impulse response $h(t) = s^*(-t)$, producing a compressed pulse with width $\tau_c \approx 1/B$ and time-bandwidth product $T_p \cdot B$, also known as the pulse compression ratio. A chirp with $T_p = 100~\mu$s and $B = 10$~MHz achieves a compression ratio of 1000, effectively providing the range resolution of a 0.1~$\mu$s pulse (15~m in range) with the energy of a 100~$\mu$s pulse.

The output of the matched filter exhibits range sidelobes that can mask weak targets adjacent to strong returns. Sidelobe reduction techniques apply amplitude weighting (windowing) to the transmitted waveform or received signal, trading peak sidelobe level for increased mainlobe width \cite{Richards2022}. Common windows include Hamming, Hanning, and Taylor weighting functions. For maritime radar applications, sidelobe levels below $-40$~dB are typically required to prevent large ship returns from obscuring nearby small targets.

The Simrad Halo~20+ radar employs pulse compression in combination with beam sharpening, a proprietary processing technique that enhances azimuthal resolution beyond the physical beamwidth of the antenna \cite{AlSharabi2021}. Details of the beam sharpening algorithm are not publicly disclosed, but the technique likely involves some form of super-resolution processing or deconvolution applied to consecutive azimuth samples. The resulting improvement in angular resolution enables better separation of closely spaced targets, which is particularly beneficial in congested harbor environments.

Pulse compression introduces additional considerations for CFAR detection. The matched filter output is no longer independent across range cells due to the correlation introduced by the sidelobe structure. Reference cells in CFAR processing must be spaced sufficiently far from the CUT to ensure statistical independence, which may require larger guard cell regions than would be necessary for a simple pulsed radar. Furthermore, the sidelobe structure can create false detections if CFAR thresholds are not designed to account for the spatially correlated nature of the compressed clutter returns \cite{Alsalem2023}.

\subsection{Simrad Halo~20+ Characteristics}

The Simrad Halo~20+ represents a consumer-grade X-band pulse-Doppler marine radar designed for recreational and light commercial vessels. Operating at X-band frequencies near 9.4~GHz, the system features a compact 20-inch radome antenna and achieves a maximum instrumented range of 36 nautical miles (67~km), though effective detection ranges depend on target size and sea conditions \cite{AlSharabi2021}. The antenna rotates at 60 revolutions per minute (RPM), corresponding to a scan rate of 1~Hz, which provides timely updates for collision avoidance while allowing sufficient dwell time for target detection.

The Halo~20+ employs pulse compression and proprietary beam sharpening algorithms to enhance both range and azimuthal resolution beyond the physical limitations of the compact antenna. The manufacturer specifies horizontal beamwidth of approximately 1.9 degrees after beam sharpening, compared to the roughly 4-degree physical beamwidth of the 20-inch antenna. This enhanced angular resolution improves target separation in congested environments such as harbors and shipping channels.

VelocityTrack Doppler processing provides target velocity information by measuring the Doppler frequency shift of moving targets relative to the ship's own motion \cite{AlSharabi2021}. The Doppler capability enables discrimination between stationary clutter (land, buoys) and moving targets (other vessels), improving detection performance in littoral environments. However, the Doppler resolution and unambiguous velocity range of the Halo~20+ are not publicly documented, limiting quantitative analysis of Doppler performance.

The radar offers four selectable operating modes tailored to different environmental conditions: Harbour mode optimizes sensitivity for short-range detection in calm waters; Offshore mode balances sensitivity and clutter rejection for open-ocean navigation; Weather mode enhances precipitation detection for storm avoidance; and Bird mode increases sensitivity for detecting flocks of seabirds. These modes likely adjust CFAR parameters, sea clutter filtering, and signal processing gain to match the expected clutter statistics and target characteristics of each scenario.

Mini automatic radar plotting aid (MARPA) functionality enables automatic tracking of 10--20 targets, computing their course and speed vectors and predicting potential collision threats \cite{AlSharabi2021}. The MARPA implementation in the Halo~20+ is not described in public literature, but typical consumer-grade systems use simplified tracking algorithms such as nearest-neighbor association and $\alpha$-$\beta$ filtering rather than the more sophisticated probabilistic data association and Kalman filtering employed in military and commercial shipping radars.

Comparison of consumer radars such as the Halo~20+ to military and research-grade systems reveals significant differences in performance, flexibility, and data access. Military radars typically provide raw or minimally processed I/Q data, enabling researchers to implement custom signal processing chains and characterize system performance in detail. Consumer radars generally output only processed PPI imagery or detection lists, with proprietary signal processing that cannot be modified or fully characterized. Detection thresholds, clutter models, and CFAR parameters are often hardcoded or adjusted automatically by undocumented algorithms, making it difficult to predict or optimize performance for specific applications such as radar-camera cross-labeling.

Research-grade radars such as those used in the IPIX, Dartmouth, and Stare datasets offer calibrated measurements with known antenna patterns, transmit waveforms, and receiver characteristics \cite{Gangeskar2022}. These systems enable controlled experiments and repeatable performance evaluations. Consumer radars lack calibration data, exhibit unknown sidelobe levels and beamwidth, and may apply non-linear post-processing (gamma correction, contrast enhancement) to the displayed imagery, complicating quantitative analysis.

Despite these limitations, consumer radars offer significant advantages for autonomous surface vehicle and small-boat applications: low cost, compact form factor, low power consumption, and simple integration with standard marine electronics. For radar-guided camera labeling, the Halo~20+ provides a practical platform for developing and demonstrating cross-sensor fusion techniques in realistic maritime environments, though researchers must carefully characterize detection performance and clutter statistics empirically.

\subsection{Track-Before-Detect Methods}

Traditional radar detection processing applies a threshold to each range-azimuth cell independently, discarding sub-threshold returns as noise or clutter. This approach is optimal when target signal-to-clutter ratio (SCR) is sufficiently high, but fails when targets produce weak returns that fall below the single-scan detection threshold. Track-before-detect (TBD) methods address this low-SCR scenario by integrating energy over multiple scans before making detection decisions, coherently accumulating weak target signatures that would be missed by conventional CFAR detectors \cite{Blackman1999}.

Particle filter TBD approaches represent the target state (position, velocity) as a weighted set of particles propagated forward in time according to a motion model \cite{Sangston2012}. At each radar scan, particle weights are updated based on the likelihood of the observed radar measurements given the hypothesized target location. Particles corresponding to true targets receive high weights through repeated correlation with target returns, while particles in clutter-only regions accumulate low weights and are eventually resampled away. Detection decisions are made by thresholding the particle weight distribution or by identifying clusters of high-weight particles.

Greco et al. developed a two-stage TBD architecture for maritime radar that combines non-coherent integration of radar magnitude measurements with recursive Bayesian filtering \cite{Greco2013}. The first stage sums radar returns across $N$ consecutive scans for each hypothesized target trajectory, producing an integrated detection statistic with improved SNR. The second stage applies a particle filter or extended Kalman filter to track promising detections and resolve ambiguities. The method demonstrated improved detection of small boats in moderate sea clutter compared to conventional CFAR processing, but at the cost of increased latency due to the multi-scan integration window.

For radar-camera cross-labeling applications, TBD methods present both opportunities and challenges. The improved low-SCR detection performance could enable labeling of distant or small targets that would be missed by single-scan CFAR. However, the multi-scan latency complicates time synchronization between radar detections and camera frames, potentially requiring retroactive association of TBD detections to camera imagery captured several scans earlier. Furthermore, TBD algorithms require accurate motion models, which may be difficult to specify for heterogeneous maritime targets exhibiting diverse maneuvering behaviors.

\subsection{DBSCAN and ST-DBSCAN for Radar Clustering}

After CFAR detection, radar returns typically consist of a collection of detections in range-azimuth space, including both true target returns and residual false alarms from clutter. Clustering algorithms group spatially proximate detections that likely originate from the same physical target, enabling target centroid estimation and facilitating downstream tracking and classification \cite{Chen2020}.

The Density-Based Spatial Clustering of Applications with Noise (DBSCAN) algorithm, introduced by Ester et al., identifies clusters as dense regions of points separated by low-density areas \cite{Ester1996}. DBSCAN requires two parameters: $\varepsilon$, the maximum distance between two points to be considered neighbors, and MinPts, the minimum number of points required to form a dense region. A point is classified as a core point if at least MinPts points (including itself) lie within distance $\varepsilon$. Clusters are formed by connecting core points that are neighbors, and non-core points that are neighbors of core points are assigned to the same cluster. Points that are not core points and have no core neighbors are labeled as noise.

For radar target clustering, $\varepsilon$ is typically set based on the expected spatial extent of target returns (a few range bins and azimuth cells), while MinPts is chosen based on the minimum number of detections expected from a point target or the minimum extent of distributed targets of interest \cite{Wang2022}. DBSCAN effectively separates closely spaced targets and rejects isolated false alarms, but it operates on a single radar scan and does not exploit temporal correlation across scans.

The Spatio-Temporal DBSCAN (ST-DBSCAN) algorithm extends DBSCAN to three-dimensional spatio-temporal space by defining separate distance thresholds for spatial and temporal dimensions \cite{Birant2007}. Two detections are considered neighbors if their spatial separation is less than $\varepsilon_1$ and their temporal separation is less than $\varepsilon_2$. ST-DBSCAN forms clusters of detections that are persistent over multiple scans, providing a simple form of track-before-detect that improves detection of low-SCR targets and rejects transient false alarms.

Zhou et al. applied ST-DBSCAN to maritime radar data for automatic identification of ship trajectories, demonstrating improved tracking continuity compared to single-scan DBSCAN followed by frame-to-frame association \cite{Zhou2024}. The spatio-temporal clustering naturally handles temporary detection dropouts caused by sea clutter or multipath fading, as detections from adjacent scans bridge the gaps. However, ST-DBSCAN assumes that cluster membership is transitive across time, which may fail for crossing targets or scenarios where two targets temporarily merge and then separate.

For radar-guided camera labeling, clustering algorithms such as DBSCAN and ST-DBSCAN provide essential preprocessing to consolidate multiple detections from extended targets into single representative points that can be associated with camera bounding boxes. The choice of clustering parameters affects the granularity of radar detections: conservative settings (small $\varepsilon$, large MinPts) produce fewer, higher-confidence clusters but may split single targets into multiple fragments, while liberal settings (large $\varepsilon$, small MinPts) merge neighboring targets and admit more false alarms. Empirical tuning based on the specific radar characteristics and target types is necessary to achieve optimal performance.

\subsection{Maritime Detection Challenges}

X-band marine radar detection performance is degraded by a variety of environmental and system-level phenomena that complicate signal processing and limit achievable detection ranges and accuracies. Understanding these challenges is essential for designing robust radar-camera cross-labeling systems and for interpreting the limitations of radar-derived labels.

Sea clutter arises from backscatter by wind-driven ocean surface waves and represents the dominant source of false alarms in maritime radar \cite{Ward2006}. Clutter power increases with sea state, radar frequency, and grazing angle. X-band radars are particularly susceptible to sea clutter due to the short wavelength, which produces strong resonant backscatter from centimeter-scale capillary waves. Sea spikes, intermittent high-amplitude clutter returns caused by breaking waves or whitecaps, create heavy-tailed clutter statistics that violate the Rayleigh assumption and necessitate more sophisticated clutter models such as the K-distribution or Pareto distribution \cite{Watts2022}.

Multipath propagation occurs when radar signals reflect from the water surface before reaching the target or after being scattered by the target, creating multiple signal paths with different delays and Doppler shifts \cite{Rosenberg2013}. Constructive interference between direct and multipath returns produces lobing patterns in the elevation plane that cause target amplitude to fluctuate as the target range or ship motion changes the multipath geometry. Destructive interference creates nulls where targets may disappear entirely. Multipath effects are most severe for low-elevation targets at moderate ranges (1--10~km) and calm sea states, where the water surface acts as a specular reflector.

Rain clutter results from backscatter by precipitation, with clutter power proportional to rainfall rate and radar frequency \cite{Skolnik2008}. X-band radars are more vulnerable to rain clutter than lower-frequency S-band or L-band systems because rain backscatter cross-section increases approximately with the fourth power of frequency (Rayleigh scattering regime). Heavy rain can produce clutter returns that obscure targets at ranges beyond a few kilometers, effectively blinding the radar. Polarization diversity techniques, which exploit differences in the polarization signatures of rain and targets, can mitigate rain clutter but are not available on most consumer marine radars \cite{Humphreys2022}.

Electromagnetic interference from other radars, communication systems, and shipboard electronics corrupts radar measurements and creates false detections. Harbor environments are particularly challenging due to the density of radar-equipped vessels and shore-based surveillance systems. Interference appears as streaks or speckle in the radar PPI display, depending on whether the interfering signal is synchronous or asynchronous with the victim radar. Frequency agility, pulse-to-pulse frequency hopping, and interference blanking circuits mitigate some forms of interference, but consumer radars offer limited protection compared to military systems \cite{DHS2009}.

Small target detection remains a fundamental challenge for X-band marine radar due to the weak radar cross-section of objects such as kayaks, debris, and buoys, which may be smaller than 1~square meter \cite{Roldan2024}. At ranges beyond a few hundred meters, such targets produce returns that are comparable to or weaker than sea clutter, resulting in low SCR and high miss rates. Increasing transmit power or antenna gain improves detection range but is constrained by regulatory limits, power consumption, and platform size. Coherent integration over multiple pulses or scans improves SCR but requires accurate motion compensation and is limited by target decorrelation time.

For autonomous vessels and radar-camera cross-labeling systems operating in cluttered maritime environments, these detection challenges directly impact the quality and completeness of radar-derived camera labels. Periods of heavy rain or high sea states will produce gaps in labeled datasets unless supplemented by other sensors. Small or low-RCS targets may be systematically underrepresented in training data, biasing learned models toward larger, higher-SNR objects. Multipath-induced amplitude fluctuations can cause intermittent detection dropouts, fragmenting trajectories and complicating track-based labeling approaches. Robust system design must account for these limitations and incorporate validation mechanisms to detect and flag low-confidence labels.

\subsection{Summary and Research Gaps}

This chapter has reviewed the fundamental signal processing techniques underlying X-band marine radar detection, including CFAR algorithms, sea clutter statistical models, pulse compression, clustering methods, and track-before-detect approaches. The performance of radar-guided camera labeling systems depends critically on the reliability of these processing stages, as errors or limitations in radar detections directly propagate to errors in automated camera labels. Classical CFAR algorithms such as CA-CFAR, OS-CFAR, and SO-CFAR provide computationally efficient adaptive detection with well-characterized false alarm rates, while emerging deep learning CFAR methods promise improved performance in complex clutter environments at the cost of increased computational demand and training data requirements. Sea clutter modeling using K-distribution, Pareto, or compound Gaussian frameworks enables accurate prediction of detection performance and informs CFAR parameter selection for heterogeneous maritime environments.

The Simrad Halo~20+ radar, representative of consumer-grade marine radar systems, offers a practical platform for developing autonomous vessel sensing and cross-sensor labeling techniques. However, significant gaps exist in published performance characterization of consumer radars compared to research-grade systems. Detailed measurements of Halo~20+ detection probability, false alarm rate, range accuracy, and angular resolution under controlled conditions are not available in the open literature, complicating performance prediction and system design. Future work should conduct systematic empirical evaluations of consumer marine radars across varying sea states, target types, and ranges to establish baseline performance metrics.

The relationship between consumer radar characteristics and research-grade radar performance represents another important gap. While expensive research radars provide superior calibration, data access, and flexibility, the extent to which their performance advantages translate to improved detection outcomes for practical autonomous vessel applications remains unclear. Comparative studies deploying both consumer and research radars on the same platform under identical conditions would quantify the performance-cost tradeoffs and inform sensor selection for resource-constrained autonomous systems.

Automation of radar-camera cross-labeling, particularly under challenging environmental conditions such as high sea states, rain, or multipath propagation, requires further research. Existing fusion and labeling approaches generally assume reliable radar detections, but as discussed in this chapter, detection performance degrades significantly in adverse conditions. Adaptive labeling confidence estimation that accounts for environmental factors and radar signal processing limitations would enable more robust dataset curation by flagging or rejecting low-quality labels.

Real-time implementation of advanced signal processing techniques such as deep learning CFAR, particle filter TBD, and ST-DBSCAN on embedded platforms typical of small autonomous vessels (e.g., NVIDIA Jetson, Raspberry Pi) presents computational challenges. Consumer radars such as the Halo~20+ output processed detections rather than raw I/Q data, limiting the applicability of sophisticated signal processing algorithms that require access to the full radar signal. Future radar designs for autonomous vessel applications should balance processing capability, cost, and power consumption while providing sufficient data access to enable custom processing pipelines.

In summary, X-band marine radar signal processing techniques provide the foundation for radar-guided camera labeling systems, but significant research opportunities remain in characterizing consumer radar performance, developing robust labeling methods for degraded detection conditions, and enabling real-time advanced signal processing on resource-constrained platforms. Addressing these gaps will enhance the reliability and scalability of automated maritime dataset construction for autonomous vessel perception.
